\section{Empirical testing of Domar aggregation}
\label{appendix:validation_Domar}

In this Appendix, we test the extent to which Domar aggregation ``works'' empirically, that is, whether Domar-aggregated industry-level TFP growth matches with aggregate TFP growth. 

We do this at various horizons, that is, we test the case that is relevant for us: assuming Domar weights are fixed for year $t$, and we are given industry level growth rates of TFP for a long period (say 20 years), will Domar aggregation reproduce the long period growth rate of aggregate TFP?

We use the EU-KLEMS data \citep{bontadiniEUKLEMSINTANProdIndustry2023}, which allows us to look at 23 countries and 56 sectors, over time spans of 10-25 years, depending on the country. We use the ``Growth Accounts Basic'' and the ``National Accounts'' files of the ``Statistical Module''. 



We compare our results to the (``second-order'') aggregation formula from \citet{baqaeeMacroeconomicImpactMicroeconomic2019}.

\subsection{Setup}

Recall that Domar aggregation is the assumption that, for a country $c$,
\begin{equation}
    \hat{A}_{c} = \underbrace{\sum_{j=1}^N \frac{p_jC_j}{pY}\hat{C}_j}_{\text{real GVA growth}}-
    \underbrace{\sum_{f=1}^F \frac{w_fL_f}{pY}\hat{L}_f }_{\text{factor quantity growth}}= 
 \sum_{i=1}^N \underbrace{\frac{p_iX_i}{pY}}_{\text{Domar weights}} \left( \underbrace{    \hat{X}_i-
    \sum_{j=1}^N\frac{p_jX_{ij}}{p_iX_i}\hat{X}_{ij}-
    \sum_{f=1}^F \frac{w_fL_{fi}}{p_iX_i}\hat{L}_{fi}
}_{\text{Industry TFP growth}}
    \right).
\end{equation}

We first describe how we compute year-on-year variables, and then how we translate them into multiple-horizon variables (i.e. ``long'' growth rates)


We get the LHS directly from EU-KLEMS\footnote{
%
KLEMS uses Tornqvist weights from factor and input shares, which introduces a slight internal inconsistency with our framework.
%
},
\[
\hat{A}_{ct} = \text{LP1ConTFP}_c
\]
where the subscript $c$ means that we read this variable for country $c$ and ``industry'' TOT. 

For the RHS, we compute Domar weight of industry $j$ in country $c$ as
\begin{equation}
    \lambda_{jct} = \frac{p_{jct} X_{jct}}{p_{ct} Y_{ct}} = \frac{\text{GO\_CP}_{jct} }{\text{VA\_CP}_{ct} },
    \label{eq:lambdajct}
\end{equation}
and industry-level TFP growth as
\[
\hat{A}_{jct} = \text{LP1ConTFP}_{jct}.
\]

We compute long growth rates by constructing aggregate TFP levels, in base 100 where the first year is equal to 100, that is 
\[
A_{j,c,t} = 100 \times \prod_{t'={t_0}}^t  \exp{(\hat{A}_{j,c,t'})},
\]
where $t_0$ is the first year where a growth rate is available in the data. We do the same for the aggregate TFP $A_{c,t}$. Thus, the ``long'' growth rate of TFP between $t$ and $t+\tau$ is
\begin{equation}
    \hat{A}_{j,c,t,\tau} \equiv \ln(A_{j,c,\tau})-\ln(A_{j,c,t})
    \label{eq:hatAjct}
\end{equation}
at the industry level, and
\begin{equation}
    \hat{A}_{c,t,\tau} \equiv \ln(A_{c,\tau})-\ln(A_{c,t})
    \label{eq:hatAct}
\end{equation}
in the aggregate.

Our Domar aggregation-based prediction for country $c$'s TFP growth rate between $t$ and $t+\tau$ is
\begin{equation}
 \hat{A}_{c,t,\tau}^\text{pred} = \sum_{j=1}^N \lambda_{j,c,t} \hat{A}_{j,c,t,\tau},
\label{eq:Acttau_pred}
\end{equation}
with elements on the RHS taken from Eqs. \ref{eq:lambdajct} and \ref{eq:hatAjct}. This can be compared to the ``true'' aggregate TFP growth, retrieved from Eq. \ref{eq:hatAct}. We measure the error $\varepsilon$, for country $c$, between $t$ and $t+\tau$ as the difference between the two,
\begin{equation}
    \varepsilon_{c,t,\tau}=  \hat{A}_{c,t,\tau}^\text{pred} - \hat{A}_{c,t,\tau}.
    \label{eq:varepsiloncttau}
\end{equation}

\subsection{Results}

We first show results at the country level, then conclude with our pooled estimates. 


\paragraph{Country-level results.} 


\begin{figure}[H]
    \centering
    \includegraphics[width=\textwidth]{figures/hindcasting_hulten/g_cumsum_error_hulten_country_secondorder.jpeg}
    \caption{Aggregation error using Hulten's theorem.}
    \label{fig_cumsum_yearly_error_hulten_country}
\end{figure}

Before proceeding to aggregating errors to obtain estimates of uncertainty at each horizon, we fix the forecast origin at the first year available in each dataset, and compute the predicted and actual TFP growth between $t$ and $t+\tau$, for each country. That is, we plot the growth rate for each horizon starting from the earlier date possible $t_0$, Eq. \ref{eq:Acttau_pred} with $t=t_0$.\footnote{
%
Note that an alternative, defining cumulative productivity as the sum of yearly growth rates, would imply that Domar weights are updated each year. We explicitly want to avoid this because in practice we will make predictions for 2050 assuming Domar weights from today.
%
}
Fig. \ref{fig_cumsum_yearly_error_hulten_country} shows the results. For many countries, the results are visually very good, but for other we observe large errors. 

While \citet{hultenGrowthAccountingIntermediate1978} justified Domar aggregation in a general equilibrium setting, \citet{baqaeeMacroeconomicImpactMicroeconomic2019} showed that under more general conditions Hulten's theorem only hold to first order. To second order, one needs to account for changes to the Domar weights, and use
\[
 \hat{A}^{\text{second order}}_{c,t,\tau}=\sum_{j=1}^N \frac{\lambda_{j,t}+\lambda_{j,\tau}}{2} \hat{A}_{j,c,t,\tau}   .
\]
While it tends to be a bit better, it does not make a big difference\footnote{
%
\cite{baqaeeMacroeconomicImpactMicroeconomic2019} find impressive difference from 1950 and 2020, and implicitly show (Fig. 8, p1193) that there is little difference between the first order Hulten, second order Hutlen and ``non linear'' (i.e. aggregate) TFP for the first 30-40 years. 
%
}; in particular, it does not resolve the large gap observed in some countries. 



Our goal here is to compute some measure of error associated with Domar aggregation. To do this, we start by aggregating different forecast origins, for a given country and forecast horizon,
\begin{equation}
\label{eq_mean_error_country_hulten}
    \bar{\varepsilon}_{c,\tau}=  \frac{1}{N_{c,\tau}}\sum_t \varepsilon_{c,t,\tau},
\end{equation}
using $\varepsilon_{i,t+\tau}$ from Eq. \ref{eq:varepsiloncttau}. Fig. \ref{fig_error_hulten_country_UK} shows the results. The errors are quite heterogeneous between countries. At horizon 11 years (where half the time series end) the mean error is about 100bp (see \ref{fig_error_hulten_quantiles}), ranging between  +1040bp (Lituania) and -600bp (Poland). The minimum error at horizon 11Y is Italy, with 1.2bp. Domar aggregation in the UK performs particularly poorly, due to errors in the early years that do not resolve in the later years.

The key conclusion from this heterogeneity is that we should calibrate the error of Domar aggregation on a larger set of countries than the UK.

\begin{figure}[H]
    \centering
    \includegraphics[width=\textwidth]{figures/hindcasting_hulten/g_error_hulten_country_UK.jpeg}
    \caption{Mean error of Domar aggregation at various forecast horizons, by country.}
    \label{fig_error_hulten_country_UK}
\end{figure}



\paragraph{Pooled estimates of the errors}

Due to the heterogeneity between countries, to get a more robust estimate of the Domar aggregation errors, we aggregate the errors across countries. For each horizon $\tau$, we simply average all the horizon-$\tau$ predictions,
\begin{equation}
\label{eq_mean_error_hulten}
    \bar{\varepsilon}_{\tau}= 
    \frac{1}{N_c}\sum_{c} \frac{1}{N_{c,\tau}}\sum_t\varepsilon_{c,t,\tau} = \frac{1}{N}\sum_{c,t} \varepsilon_{c,t,\tau},
\end{equation}
where $N_c$ the number of countries, and $N_{c,\tau}$ the length of the country $c$ time series minus $\tau$ ($N$ is the number of observations crossing countries and pairs of years separated by $\tau$ years).

\begin{figure}[H]
    \centering
    \includegraphics[width=\textwidth]{figures/hindcasting_hulten/g_error_hulten_quantiles.jpeg}
    \caption{Mean, median and quantiles of error of Hulten theorem}
    \label{fig_error_hulten_quantiles}
\end{figure}


Fig. \ref{fig_error_hulten_quantiles} shows the mean error along with the quantiles at 10, 25, 50, 75 and 90\% across all industries and countries for each horizon $\tau$ (Eq. \ref{eq_mean_error_hulten}). The bias is remarkably low, even at long horizon, but with relatively important errors around the mean. This shows that as long as we take uncertainty into account, using Domar aggregation is perfectly justified.


\paragraph{Robustness.} We intend to repeat this test at a finer level of aggregation, and on a different database, OECD STAN.



